## TEMP holding note — RQ1's own XGBoost hyperparameter search (2026-08-16)

**Context**: RQ1's "fairly tuned" XGBoost row (`TEMP_rq1_baseline_comparison_2026-08-11.tex`) used
a config (`n_estimators=500, max_depth=4, learning_rate=0.04`) borrowed from RQ3's
`explain_signal.py`, itself never validated by a real search -- reused across RQ1/RQ2b/RQ3 as a
"borrowed, not searched" default. This note replaces that hedge for RQ1 specifically with a real
per-target search: **only RQ1's XGBoost was retuned. RQ2b and RQ3's XGBoost configs are
deliberately left untouched** -- retuning RQ3 in particular carries real blast-radius risk, since
its entire outlier-diagnosis thread (which specific plots are worst-residual) depends on the
current fit.

**Method**: `models/baselines/rq1_xgb_hyperparameter_search.py`. Grid: `n_estimators` in
{300, 500, 800}, `max_depth` in {3, 4, 6}, `learning_rate` in {0.02, 0.04, 0.08} = 27 configs,
each fit on Set3 (`nested_set3_gated_terrain_wind_vif`), both cohorts, `spatial_block_kfold`
(5 folds) = 270 fits total. **Model selection uses validation R2 only, never test R2**, to avoid
tuning-on-test leakage; test R2/RMSE/MAE reported only for the winning config, after selection is
already decided. Reuses the same data/split pipeline as
`xgb_hyperparameter_sensitivity_check.py`'s `fit_and_score_rq1_fold()`.

### Winning config per cohort (selected on mean validation R2 across 5 folds)

| Cohort | n_estimators | max_depth | learning_rate | Val R2 | Test R2 | Test RMSE | Test MAE |
|---|---:|---:|---:|---:|---:|---:|---:|
| 4survey | 500 | 6 | 0.02 | 0.682 | 0.674±0.009 | 4.548±0.137 | 3.431±0.123 |
| 6survey | 300 | 3 | 0.08 | 0.717 | 0.722±0.031 | 3.656±0.232 | 2.669±0.183 |

### Borrowed config's result, for direct comparison (same pipeline, same folds)

| Cohort | Val R2 | Test R2 | Test RMSE | Test MAE |
|---|---:|---:|---:|---:|
| 4survey | 0.678 | 0.675±0.009 | 4.539±0.152 | 3.431±0.153 |
| 6survey | 0.712 | 0.718±0.025 | 3.682±0.221 | 2.697±0.170 |

**Finding: a real 270-fit search barely moves the numbers.** The winning config's test R2 sits
within 0.002-0.004 of the borrowed config's on both cohorts -- well inside the fold-to-fold SD
(0.009 on 4survey, 0.025-0.031 on 6survey). RMSE/MAE are likewise indistinguishable (4survey MAE
is identical to three decimal places: 3.431 vs. 3.431). On 4survey the winning-by-val-R2 config is
actually marginally *lower* on test R2 than the borrowed config (0.674 vs. 0.675) -- expected
noise from selecting on validation rather than test, not a sign the search failed.

**What this means for RQ1 item 1**: the "fairly tuned" hedge (borrowed, unsearched config) can now
be dropped -- a genuine per-target search was run and it confirms the borrowed config was already
close to optimal for this task. This also strengthens (does not weaken) the headline claim that
XGBoost beats every neural model: the margin isn't an artifact of a lucky or unfair hyperparameter
choice, since 270 alternative configs, honestly selected on validation data, don't find anything
meaningfully better.
